% ════════════════════════════════════════════
% 模型检验与分析（6.模型检验.tex）写作规则 —— 模板内嵌，AI 先读本文件再写
% ════════════════════════════════════════════
\section{模型检验与分析}
\label{sec:validation}

我们从预测误差和条件变化两方面检验模型：前者考察对检验记录的预测，后者考察赋权、成本与预算改变后配置及结论是否保持。

\subsection{误差分析}
\label{sec:check}

表\ref{tab:error}汇总各问使用的检验指标。问题一的线性岭与核岭、问题二的经典和广义标度律，以及问题四的损失--得分桥接和合并横截面回归，分别报告 $R^2$、RMSE、MAE、MAPE。我们据此同时评价拟合解释度与预测偏差，避免仅凭相关性判断效果。

设同一评价集合中的观测值为 $y_i$，预测值为 $\hat y_i$，记录数为 $n$，均值为 $\bar y$。几类指标分别由残差的平方、绝对值及相对大小构成：
\begin{align*}
 R^2&=1-\frac{\sum_{i=1}^{n}(y_i-\hat y_i)^2}{\sum_{i=1}^{n}(y_i-\bar y)^2},\\
 \mathrm{RMSE}&=\sqrt{\frac{1}{n}\sum_{i=1}^{n}(y_i-\hat y_i)^2},\qquad
 \mathrm{MAE}=\frac{1}{n}\sum_{i=1}^{n}|y_i-\hat y_i|,\\
 \mathrm{MAPE}&=\frac{100\%}{n}\sum_{i=1}^{n}\left|\frac{y_i-\hat y_i}{y_i}\right|.
\end{align*}
其中，$R^2$ 需要观测值具有非零变异，MAPE 的逐条计算则要求观测值非零。RMSE 对较大的偏差给予更多权重，MAE 描述平均绝对偏离，两者与被预测量具有相同量纲；MAPE 便于阅读相对误差，但观测值靠近零时，少量记录也可能显著影响其大小。因此，我们结合样本取值范围共同报告这些指标。

我们还区分决定系数与相关系数的用途。$R^2$ 将残差与观测值围绕均值的变异作比较，要求预测值在实际尺度上接近观测值；Pearson 的 $r$ 则衡量两者共同变化的程度。预测序列整体偏高或偏低时，仍可能保持较高相关。表中的跨尺度相关用于评价趋势，其余结果按各问的评价集合解释，不据同列数值直接排列所有模型的优劣。

\begingroup
\setlength{\tabcolsep}{3pt}
\begin{longtable}{@{}>{\raggedright\arraybackslash}p{0.31\textwidth}>{\centering\arraybackslash}p{0.17\textwidth}>{\centering\arraybackslash}p{0.15\textwidth}>{\centering\arraybackslash}p{0.13\textwidth}>{\centering\arraybackslash}p{0.13\textwidth}@{}}
\caption{各问拟合、预测误差与预算约束残差汇总}\label{tab:error}\\
\toprule
  模型（数据集） & $R^2$ / $r$ 或残差 & RMSE / 相对残差 & MAE & MAPE(\%) \\
\midrule
\endfirsthead
\caption[]{各问拟合、预测误差与预算约束残差汇总（续）}\\
\toprule
  模型（数据集） & $R^2$ / $r$ 或残差 & RMSE / 相对残差 & MAE & MAPE(\%) \\
\midrule
\endhead
\bottomrule
\endfoot
问题一岭回归（训练 1M，$n=512$） & $R^2=0.629$ & 0.551 & 0.403 & 8.54 \\
  问题一岭回归（检验 1M，$n=256$） & $R^2=0.619$ & 0.523 & 0.386 & 8.43 \\
  问题一核岭（检验 1M，$n=256$） & $R^2=0.9454$ & 0.2256 & 0.1344 & — \\
  问题一跨尺度 Pearson（60M / 1B） & $r$: 0.782 / 0.651 & — & — & — \\
  问题二经典标度律（B1，$n=1{,}176$） & $R^2=0.9999998$ & $1.5\times10^{-4}$ & $1.1\times10^{-4}$ & 0.004 \\
  问题二经典标度律留出（4 小模型$\to$4 大模型） & $R^2=0.9999997$ & $1.5\times10^{-4}$ & — & 0.005 \\
  问题二广义标度律（B1+B6 联合，$n=1{,}536$） & $R^2=0.990$ & 0.036 & 0.017 & 0.66 \\
  问题二跨族验证（B4，$n=57$） & $r=0.975$ & 0.293 & — & 8.83 \\
  问题二文献验证（B5，$n=44$） & $r=0.910$ & 0.193 & — & 7.99 \\
  问题四 Loss—Benchmark 桥接（C6，$n=75$） & $R^2=0.26$ & — & — & 67.79 \\
  问题四合并横截面回归（C1 开源筛选，$n=2{,}346$） & $R^2=0.475$ & — & — & — \\
  问题四合并横截面回归（C3 长周期，$n=4{,}579$） & $R^2=0.424$ & — & — & — \\
  问题四前沿 logistic（历史拟合） & $R^2=0.975$ & — & — & — \\
  问题三优化约束残差（$|C_{\text{total}}/C_{\text{budget}}-1|$） & — & $<10^{-6}$ & — & — \\
\end{longtable}
\endgroup

问题一的同尺度预测与训练表现接近。表\ref{tab:error}中，13 个验证域的平均训练 $R^2$ 为 $0.629$，1M 检验为 $0.619$，相差 $0.010$，MAPE 约 $8.5\%$。这一对照没有显示明显的训练与检验落差，但不能据此排除线性欠拟合。核岭的同尺度平均域 $R^2=0.9454$、综合 RMSE 为 $0.2256$；相应训练五折误差为 $0.2558$。跨尺度绝对误差见表\ref{tab:p1_nonlinear}，较高趋势相关不代表绝对损失准确。

问题二的精度取决于检验来源。对 B1，经典拟合有 $R^2\approx1.000$、MAPE $0.004\%$，仅用小模型估计后，对大模型的留出预测仍有 $R^2=0.9999997$；残差标准差 $1.5\times10^{-4}$ 相比损失标准差 $0.342$ 很小。这些结果说明对 B1 的贴合程度高，尚不能说明跨来源同样准确。B4 和 B5 的 MAPE 分别为 $8.8\%$、$8.0\%$，相关分别为 $r=0.975$、$r=0.910$。广义模型的 $1{,}536$ 条联合记录给出 $R^2=0.990$、MAPE $0.66\%$；由于样本组成改变，不能将其拟合优度与经典模型直接作差来判断泛化收益。

问题四的换算精度主要受样本可比性限制：中可比层的验证集不同，高可比层的得分跨度有限，合并后的 $R^2=0.26$、MAPE 为 $67.8\%$，故结果仅作方向解释。问题三检查的是预算可行性；表中最后一行的数值是成本回代得到的相对预算残差，虽与 RMSE 同列展示，不能解释为预测误差。有效解的残差均 $<10^{-6}$，符合成本约束。

各问的评价集合不同，也决定了误差数字的使用范围。问题一的检验针对新配方下的领域损失，问题二的对照涉及不同规模、质量和资料来源，问题四的桥接还要跨越损失与能力两种指标。较小误差能够支持相应集合内的描述，却不能自动扩大到其他来源。因此，我们按记录来源及其是否参与估计，限定表\ref{tab:error}中各项误差的解释范围。

预算约束残差与上述预测误差的含义不同。残差接近零，说明返回的参数量、数据量和质量在给定成本式下基本用尽预算；它本身不衡量标度律的预测精度，也不说明外部评测得分会提高多少。因此，我们从预算可行性、预测损失和资源份额三个方面评价问题三的配置。

\subsection{灵敏度分析}
\label{sec:sensitivity}

表\ref{tab:sensitivity}比较赋权方式、冲突阈值、岭正则、质量成本、质量基线、上下文长度和预算变化后的结果，同时补充指标方向与跨来源验证的影响。我们通过这些对照识别较稳定的结论，以及依赖特定成本或样本的结论。

\begin{longtable}{@{}>{\raggedright\arraybackslash}p{0.30\textwidth}>{\raggedright\arraybackslash}p{0.36\textwidth}>{\raggedright\arraybackslash}p{0.26\textwidth}@{}}
\caption{评分规则、成本与预算变化下的结果稳定性}\label{tab:sensitivity} \\
\toprule
参数/设定变化 & 输出变化 & 灵敏度评估 \\
\midrule
\endfirsthead
\caption[]{评分规则、成本与预算变化下的结果稳定性（续）} \\
\toprule
参数/设定变化 & 输出变化 & 灵敏度评估 \\
\midrule
\endhead
\midrule
\multicolumn{3}{r}{续下页} \\
\endfoot
\bottomrule
\endlastfoot
赋权方案：熵权法 $\to$ 等权平均 $\to$ 首主成分非负载荷加权 & 域级质量排序一致性 $0.643\to0.679$；样本级得分相关 $0.875/0.927$；得分均值 $0.491\to0.637\to0.496$ & 中等敏感：整体格局稳定，但领域名次会变化 \\
  指标方向：ad\_en 按“索引 1=无广告”（正向）$\to$ 若误判为负向 & 领域质量排序由「arxiv $>$ stackexchange $>$ c4 $>$ commoncrawl $>$ book $>$ wikipedia $>$ github」变为「c4 $>$ arxiv $>$ commoncrawl $>$ stackexchange $>$ book $>$ wikipedia $>$ github」 & 高度敏感：方向判定必须用实证证据 \\
  冲突阈值 $\tau$：$0.5\to1.0\to1.5\to2.0\to2.5$ & 冲突率 $0.657\to0.394\to0.219\to0.120\to0.066$；与独立基准之比 $0.91\to0.82\to0.76\to0.76\to0.86$ & 结论稳健：所考察阈值下冲突率均低于独立基准 \\
  岭正则系数 $\lambda$：$10^{-4}\sim10^{-1}$ & 训练 $R^2$ 由 $0.6291$ 降至 $0.6238$，检验 $R^2$ 由 $0.6187$ 降至 $0.6157$ & 拟合精度较稳定；系数与候选配比仍会变化 \\
  质量成本函数：指数型 $\to$ 幂函数型 $\to$ 对数渐进型（$C=10^{19}$） & $Q^{*}=0.660\to0.549\to1.000$，$L^{*}$ 差异 $<0.06$；质量份额 $13.1\%\to0.0\%\to31.7\%$ & 低预算档较敏感 \\
  同上，考察 $C=10^{22}$、$10^{24}$ & 三种函数下 $Q^{*}$ 均为 $1.000$；在 $10^{22}$ 档，$D^{*}$ 为 229.4--274.1B tokens，最大值较最小值高约 $19.5\%$ & 质量均达到上界，规模配置仍有差异 \\
  质量基线口径：配比加权 $Q_0=0.5487\to$ 样本均值 $0.4912$ & 三档预算下 $Q^{*}$：$0.549/1.000/1.000\to0.507/1.000/1.000$，低预算质量份额由 $0$ 变为 $2.63\%$，$N^{*}$ 最大变化约 $3.24\%$ & 规模配置变化较小；低预算是否启动质量投入依赖基线口径 \\
  上下文长度：$2{,}048\to131{,}072$ tokens（C7 可行取值） & 注意力/训练倍数 $0.07\to4.37$；$N^{*}$ 由 $6.091$B 降至 $2.705$B；损失上升 $6.0\%$ & 对配置影响较大；$L_{\text{crit}}=30{,}000$ 为两项成本相等的位置 \\
  算力预算 $C$：$10^{19}\to10^{24}$ FLOPs & $Q^{*}$ 由 $0.549$ 升至 $1.000$；$D^*/N^*$ 由 $30.9$ 升至 $76.1$；质量份额 $0\to0.407\to0.016$ & 本成本与质量曲率下出现三个投入阶段 \\
  半合成数据验证：B2 / B8 & B2 的 MAPE 为 $33.7\%$；B8 校准段翻转质量方向后为 $133.4\%$，其部分数据损失低于本模型的参考渐近损失基准 & 跨来源偏差较大，两者均未参与主拟合 \\
\end{longtable}

表\ref{tab:sensitivity}显示，不同设定对结果的影响程度并不相同。

我们按条件作用的位置解释这些差异。调整权重或指标方向，会先改变质量评分；调整质量费用和上下文长度，会改变预算中各项支出的比例；扩大总预算，则是在同一成本形式下改变可配置的资源总量。因而，我们除比较最终损失外，也比较中间变量和配置份额。若两种方案的预测损失接近，但质量份额相差较大，就说明目标值相近并不要求投入方式相同。

阈值、正则参数和预算扫描均由有限取值构成。相邻取值下的结果可用来判断已检验范围内的变化方向，首次发生变化的网格点则给出定位边界的线索。我们保留扫描间距和求解状态，并结合费用关系解释变化。图中连接相邻点仅用于展示趋势，未检验区间不具有与扫描点相同的数值依据。

\textbf{（1）较稳定的结果。} 在所检验阈值内，冲突率始终低于独立参照，比值均 $<0.92$。岭正则变化对拟合精度影响较小；更换质量基线后，规模配置变化较小，但低预算是否启动质量投入随基线口径改变。三类成本在所考察的中高预算档都使质量达到上界，但参数量、数据量的具体组合仍然不同。

\textbf{（2）需要区分的边界。} $L_{\text{crit}}=30{,}000$ tokens 只表示注意力费用与基础训练费用相等，不代表最优解在此处发生突变；$1.58\times10^{19}$、$7.08\times10^{19}$ FLOPs 是预算网格中首次出现质量投入和首次达到上界的有效点，实际边界还受网格间距、求解状态及成本形式影响。

\textbf{（3）依赖条件的结果。} $C=10^{19}$ FLOPs 下更换质量成本时，最优质量覆盖 $0.549\sim1.000$，说明低预算策略依赖成本设定，应同时比较各情景；损失--得分桥接只有 $R^2=0.26$，不能精确换算能力分；领域排序也会受指标方向和赋权方式影响。

补充检验进一步区分了拟合精度与决策稳定性。问题一在同一个加权损失目标下比较线性岭与核岭，见表\ref{tab:q1_review_target}及表\ref{tab:q1_review_surrogates}，发现原线性候选的相对收益方向随响应函数而变；正则变化对平均 $R^2$ 的影响虽小，系数符号和领域份额仍会变化，见表\ref{tab:q1_review_ridge}。问题二、三的表\ref{tab:q23_review_quality_scale}和表\ref{tab:q23_review_gamma}分别检查质量映射与曲率：保持质量排序不变仍会改变边际收益，而将 $\gamma$ 固定为 1 时，拟合目标只增加约 $0.024\%$，原首次饱和预算下的最优质量降至约 $0.964$。因此，精度相近不足以证明端点配置稳定。

表\ref{tab:q23_review_m_amplitude}则将配比修正的迁移幅度从零逐步提高到原值。等效预算倍数随之从 1 增至约 29.6，说明该量主要反映配比收益的迁移假设；它不能单独作为真实训练节省的估计。以上对照均另存为诊断输出，原拟合参数和正式配置结果保持不变。

问题四的表\ref{tab:q4_review_time}显示，年份组差对同名去重较稳定，但含义仍是不同模型记录之间的条件差异。表\ref{tab:q4_review_window}和表\ref{tab:q4_review_holdout}进一步表明，前沿曲线的饱和位置对平台窗口敏感，所选时间留出期的误差也高于保持最后峰值的参照。这项事后检查不足以选定一个普遍更优的预测方法，但足以说明全样本拟合度不能替代时间外验证。

\subsection{各问结果的衔接与使用条件}

完成单项检验后，我们进一步核对结果在各问之间的传递。质量基线、配比修正和标度律参数对下一问的作用不同，因此分别检查其影响位置与使用条件，明确现有检验覆盖的范围。

质量基线 $Q_0$ 一方面确定问题三的质量起点，另一方面进入增量费用 $[g(Q)-g(Q_0)]_+$。所以，更换基线会同时改变质量可选区间和费用参照，不能仅把结果表中的质量数值平移后仍视为同一方案。前述两种基线的对照已经给出相应配置；我们将各自基线、成本形式及配置共同列为使用条件。问题一中缺少直接领域质量观测的部分，也需要随映射规则一并说明。

配比修正 $M(\bm p)$ 则以损失平移量进入广义模型。配比固定时，它不参与参数量、数据量和质量的求导，但仍影响报告的预测损失水平。因而，讨论规模配置与讨论配方收益时，需要明确是否沿用同一配比。问题一的有界候选在不同响应函数下收益方向不同，问题二的跨来源对照也未完成配比响应的数值标定，因此配方收益及其算力等价只作为条件情景。

标度律参数决定不同投入的边际收益，成本函数决定实现这些投入要支出多少算力。问题二的拟合误差与问题三的成本敏感性分别考察这两侧的关系。现有结果可以说明给定参数和成本下怎样配置，也可以比较已设定成本情景之间的差别；对于参数误差与成本误差同时变化后的影响，本文没有另外给出联合概率区间，因此不把这些情景差别解释成最优配置的置信范围。

能力预测又使用另一组观测。损失--得分桥接展示了两类指标的关联程度，历史前沿和合并横截面回归则直接使用评测记录。问题三中的损失下降据此解释为训练目标的改善，问题四的能力增量保留为相应评测尺度下的结果；两者均应结合各自的样本范围和桥接假设解读。
